\section{Interpretation}
\label{sec:interpretation}

\paragraph{What drives the model.} We measured permutation importance on the test set as the
drop in PR-AUC when one feature is shuffled (\Cref{fig:perm}). For the decision tree, three
features carry almost everything: shipping mode (0.164), scheduled days (0.087, which is the
same information as mode) and order hour (0.053). Every other feature scores below 0.001. For
gradient boosting the picture is the same: mode 0.283, hour 0.035, and nothing else above 0.002.
Destination, customer history, order value, discount and product category do not matter.

\begin{figure}[!htbp]
  \centering
  \includegraphics[width=0.5\linewidth]{06_permutation_importance}
  \caption{Permutation importance for the decision tree: mean drop in test PR-AUC over 5 shuffles, with error bars showing one standard deviation.}
  \label{fig:perm}
\end{figure}

\paragraph{Tree structure.} The depth-3 tree in \Cref{fig:tree} reads as a policy manual. The
first split separates Standard Class (40\% late) from the rest. On the left, the tree
isolates Same Day (\tcode{scheduled\_days} $\le 0.5$) and splits it at \tcode{hour} $\le 11.5$
into leaves that are 0\% and 100\% late. First Class becomes a 100\%-late leaf, and Second Class
a leaf that is 80\% late. On the right, the splits on target-encoded city and city $\times$ mode
history move the late share only between 39\% and 43\%. That is noise, and these leaves
disappear from the importance ranking.

\begin{figure}[!htbp]
  \centering
  \includegraphics[width=0.85\linewidth]{07_decision_tree_depth3}
  \caption{Decision tree limited to depth 3, fitted on the training period. \tcode{value} gives the shares [on time, late]; \tcode{samples} is the share of training orders reaching each node.}
  \label{fig:tree}
\end{figure}

\paragraph{Does this make supply-chain sense?} Mostly, yes, and it points at the promise rather
than the order.
\begin{itemize}
  \item \textbf{First Class} is promised 1 day and always takes 2: a service-level design error,
        not a risk to predict.
  \item \textbf{Second Class and Standard} both take 2--6 days uniformly, mean 4.0
        (\Cref{tab:requote}). Second Class is \emph{not faster}, only promised as faster, and
        its buyers are late 80\% of the time.
  \item \textbf{Same Day} behaves like a warehouse with a 12:00 cut-off, which real fulfilment
        centres use; the fix is to show the cut-off at checkout.
  \item The \textbf{absence} of destination, customer or product effects is not plausible for
        real logistics. It is consistent with the realism checks (\Cref{subsec:realism}), so we
        read it as a property of the data.
\end{itemize}
The partial-dependence plots in \Cref{app:figures} confirm this. The average predicted late
probability jumps by about 6 points at 12:00. For order value, discount and city $\times$ mode
history it moves by at most 2 points, and a feature whose permutation importance is
essentially zero cannot turn that movement into better rankings.
